% LSTM post, figure 2: a standard RNN unrolled over three steps. Its
% repeating module holds a single tanh layer over [h_{t-1}, x_t]; drawn on
% the same module geometry and colour key as lstm-chain, for comparison.
\documentclass[tikz,border=3pt]{standalone}
\input{FIGKIT/preamble.tex}
\tikzset{
  tlay/.style={rectangle, preaction={fill=paper, draw=none}, tint=fd, line width=.7pt,
    minimum width=8.4mm, minimum height=4.6mm, inner sep=0pt, font=\footnotesize},
  wire/.style={line width=.7pt, rounded corners=2mm},
  ghost/.style={draw opacity=.13, fill opacity=0, text opacity=.16,
    tlay/.style={rectangle, draw=fd, line width=.7pt, minimum width=8.4mm, minimum height=4.6mm,
      inner sep=0pt, font=\footnotesize}},
  xin/.style={op=fa, minimum size=8mm},
  hout/.style={op=fe, minimum size=8mm},
}
% module box (0,0)-(\W,\H), as in lstm-chain; the state runs at height \yS
\def\W{52}\def\H{33}\def\yS{22}\def\xX{3}\def\xU{48}\def\gap{9}
\newcommand\Rbox{\draw[tint=fc, rounded corners=2.5pt] (0,0) rectangle (\W,\H);}
\newcommand\Rnn{%
  \draw[wire, ->] (0,\yS) -- (16,\yS) -- (16,9) -- (26,9) -- (26,13.7);
  \draw[wire] (\xX,0) -- (\xX,9) -- (20,9);
  \draw[wire] (26,18.3) -- (26,\yS) -- (\W,\yS);
  \draw[wire] (\xU,\yS) -- (\xU,\H);
  \node[tlay] at (26,16) {$\tanh$};}
\begin{document}
\begin{tikzpicture}[fig, x=1mm, y=1mm]
  \foreach \k/\s in {-1/{t-1}, 0/t, 1/{t+1}} {
    \begin{scope}[xshift=\k*(\W+\gap)*1mm]
      \Rbox
      \ifnum\k=0 \Rnn \else
        \begin{scope}[ghost] \Rnn \end{scope}
        \node[font=\LARGE] at (38,14) {$A$};
      \fi
      \node[xin] (x) at (\xX,-12) {$x_{\s}$};
      \node[hout] (h) at (\xU,\H+12) {$h_{\s}$};
      \draw[wire] (x) -- (\xX,0);
      \draw[arr] (\xU,\H) -- (h);
    \end{scope}
  }
  \draw[arr] (-\gap,\yS) -- (0,\yS);
  \draw[arr] (\W,\yS) -- (\W+\gap,\yS);
  \draw[arr] (2*\W+\gap,\yS) -- (2*\W+\gap+7,\yS);
\end{tikzpicture}
\end{document}
